\documentclass[11pt,a4paper]{article}

% ============================================
% PACKAGES
% ============================================
\usepackage[utf8]{inputenc}
\usepackage[T1]{fontenc}
\usepackage{amsmath,amsthm,amssymb}
\usepackage{mathtools}
\usepackage{enumitem}
\usepackage{booktabs}
\usepackage{geometry}
\usepackage{hyperref}
\usepackage{cleveref}
\usepackage{float}
% \usepackage{microtype} % omitted for compatibility
\usepackage{xcolor}
\usepackage{stmaryrd}

\geometry{margin=1in}

\hypersetup{
    colorlinks=true,
    linkcolor=blue!70!black,
    citecolor=green!50!black,
    urlcolor=blue!70!black
}

% ============================================
% THEOREM ENVIRONMENTS
% ============================================
\theoremstyle{plain}
\newtheorem{theorem}{Theorem}[section]
\newtheorem{lemma}[theorem]{Lemma}
\newtheorem{proposition}[theorem]{Proposition}
\newtheorem{corollary}[theorem]{Corollary}

\theoremstyle{definition}
\newtheorem{definition}[theorem]{Definition}
\newtheorem{example}[theorem]{Example}
\newtheorem{axiom}[theorem]{Axiom}

\theoremstyle{remark}
\newtheorem{remark}[theorem]{Remark}
\newtheorem{observation}[theorem]{Observation}

% ============================================
% CUSTOM COMMANDS
% ============================================
\newcommand{\N}{\mathbb{N}}
\newcommand{\Z}{\mathbb{Z}}
\newcommand{\R}{\mathbb{R}}
\newcommand{\D}{\mathbb{D}}
\newcommand{\U}{\mathcal{U}}
\newcommand{\B}{\mathcal{B}}
\newcommand{\Disc}{\mathrm{Disc}}
\newcommand{\nextmod}{\bigcirc}
\newcommand{\eventually}{\Diamond}
\newcommand{\fix}{\mathrm{fix}}
\newcommand{\Stream}{\mathrm{Stream}}
\newcommand{\unfold}{\mathrm{unfold}}

% ============================================
% TITLE
% ============================================
\title{\textbf{The Algorithmic Big Bang:\\
Cosmogenesis, Inflation, and the Arrow of Time\\
at the Univalent Horizon}}

\author{Halvor S.\ Lande\\
\texttt{hsl@awc.no}}

\date{March 2026}

% ============================================
% DOCUMENT
% ============================================
\begin{document}

\maketitle

% ============================================
% ABSTRACT
% ============================================
\begin{abstract}
General Relativity fails at $t = 0$, yielding nonsensical infinite values at the initial singularity.
Physics currently lacks a mathematical vocabulary for a universe assembling its own rules.
We demonstrate that the initial singularity is not a point of infinite mass-energy density in pre-existing spacetime, but the exact computational phase transition where external metamathematical assembly (algorithmic time, $\tau$) is compiled into internal geometric dynamics (physical time, $t$).
By applying the Kinematic-Dynamic Equivalence ($\nextmod X \simeq X^{\D}$) of the Dynamical Cohesive Topos at Step~15 of the PEN Generative Sequence~\cite{pen}, we show that the Big Bang represents the ignition of autonomous guarded corecursion.
This algorithmic origin derives the forward execution asymmetry underlying thermodynamic irreversibility, yields a sharply compressed low-entropy ignition state in Zurek's algorithmic sense, and supplies the structural precondition for any subsequent inflationary phase.
The Univalent Horizon---the structural boundary at which the sequence permanently halts---is identified as a Planck-like cosmological firewall separating the pre-physical compilation epoch from the self-computing physical universe.
\end{abstract}

% ============================================
% 1. INTRODUCTION
% ============================================
\section{Introduction: Cosmogenesis as a Logical Phase Transition}
\label{sec:introduction}

\subsection{The First-Cause Paradox and Its Dissolution}

Standard cosmological models treat time as a pre-existing real line $\R$ upon which a singularity is placed at $t = 0$.
The singularity is then a pathology: a point where curvature scalars diverge, geodesics are incomplete, and the field equations cease to predict~\cite{hawking-ellis}.
This framing converts an explanatory failure into an ontological posit---the universe ``began'' in a state that the theory's own mathematics cannot describe.

The difficulty is structural, not technical.
Any framework that presupposes a continuous time manifold and then asks what happened at its boundary is asking the container to explain its own origin.
The question ``what caused the Big Bang?'' inherits the full weight of the first-cause paradox: every causal explanation requires a prior temporal context, generating an infinite regress or an unexplained brute fact.

PEN dissolves this paradox by rejecting the premise.
Time is not an ambient container in which events are arranged; it is a data structure that had to be algorithmically constructed and selected for efficiency~\cite{pen}.
The Generative Sequence builds spatial topology, cohesion, differential structure, metrics, and functional-analytic machinery across 14~discrete steps \emph{before} temporal dynamics exists.
Physical time---continuous, geometric, directional---arrives at Step~15 as the output of the Kinematic-Dynamic Equivalence ($\nextmod X \simeq X^{\D}$), not as a background assumption.

\subsection{Ontological Stance and Scope}
\label{sec:stance}

We adopt a constructive variant of ontic structural realism.
The claim is not that the universe is a simulation running on a background computer external to physics.
Rather, the claim is that the constraints governing the emergence of a causally closed formal system in PEN are structurally isomorphic to the constraints governing the emergence of spacetime and dynamics.
The compiler/runtime vocabulary marks the distinction between external metatheoretic description and internal physical execution; it does not introduce an extra substrate behind the structure itself.

Methodologically, this paper is corollary-driven.
The existence of the 15-step Generative Sequence, the audited counts $\kappa$ and $\nu$, and the passage from the selected temporal shell to the DCT are imported from~\cite{pen}.
What is new here is the cosmological interpretation of those results: the status of Step~15 as ignition, the clarification of the two clocks, and the consequences for inflation, irreversibility, and the Past Hypothesis.
The extraction of specific phenomenological sectors---including explicit $3+1$-dimensional Lorentzian geometry, the Einstein-Hilbert layer, and related low-energy signatures---is deferred to the companion analysis~\cite{gr-companion}.

\subsection{Compiler vs.\ Runtime}
\label{sec:compiler-runtime}

PEN's cosmogenesis mirrors the distinction between compilation and execution in computation theory.
The analogy is structural: it tracks the passage from externally assembled syntax to internally generated dynamics.

\begin{definition}[Compilation Epoch]
\label{def:compilation}
The \emph{compilation epoch} (Steps~1--14) is the discrete metamathematical assembly of the kinematic framework.
During this epoch, the library $\B_n$ grows monotonically under the selection dynamics (Axioms~1--5 of~\cite{pen}), accumulating type formers, higher-inductive types, modal operators, and axiomatic extensions.
No continuous internal observers, trajectories, or dynamical equations exist.
Time is measured solely by algorithmic ticks $\tau_n$ and integration latency $\Delta_n = F_n$.
\end{definition}

\begin{definition}[Runtime Ignition]
\label{def:ignition}
\emph{Runtime ignition} (Step~15) is the selection of the temporal-cohesive shell whose semantic completion is the Dynamical Cohesive Topos (DCT): it introduces temporal modalities ($\nextmod$, $\eventually$), explicit exchange laws with cohesion, and guarded coherence clauses into the library.
At this step, the universe transitions from extrinsic structural selection (maximizing $\rho$) to intrinsic dynamical evolution (integrating vector fields on $\U$).
\end{definition}

The Big Bang, on this account, is neither a physical explosion nor a mathematical singularity.
It is the exact moment at which a self-assembling formal system becomes causally closed and begins computing its own future.


% ============================================
% 2. THE PRE-PHYSICAL EPOCH
% ============================================
\section{The Pre-Physical Epoch: Algorithmic Time \texorpdfstring{($\tau$)}{(tau)}}
\label{sec:pre-physical}

\subsection{Assembly of the Phase Space}
\label{sec:assembly}

Prior to Step~15, the library $\B_{14}$ contains the full spatial and functional-analytic infrastructure of modern physics: dependent types, higher-inductive types ($S^1, S^2, S^3$), cohesion ($\flat, \sharp, \Disc, \Pi_{\mathrm{coh}}$), connections, curvature, metrics, frame bundles, and Hilbert functionals~\cite{pen}.
These structures exist as a static, well-typed formal library---a univalent 0-groupoid sequence of realized specifications.
They describe \emph{what kinds of spaces and operations are available}, not \emph{how anything moves within them}.

The distinction is categorical: $\B_{14}$ is a kinematic library, not a dynamical system.
It contains the type $S^3$ but no trajectory on $S^3$; it contains the metric tensor $g_{\mu\nu}$ but no Einstein equation governing its evolution; it contains Hilbert spaces but no Schr\"odinger equation propagating states through them.
The library is a phase space without a Hamiltonian---a stage fully constructed but with no actors yet in motion.

\subsection{The Two Clocks}
\label{sec:two-clocks}

A persistent conceptual error in cosmological thinking is the implicit assumption that ``before the Big Bang'' must refer to a prior interval of physical time.
PEN makes the temporal structure fully explicit by distinguishing two incommensurable notions of time.

\begin{definition}[Algorithmic Time]
\label{def:alg-time}
\emph{Algorithmic time} $\tau$ is the discrete tick counter of the Generative Sequence.
At each tick $\tau_n$, the selection dynamics evaluates candidates against the bar $\mathrm{Bar}_n = \Phi_n \cdot \Omega_{n-1}$, realizes the minimal-overshoot winner, and increments the library.
The elapsed algorithmic time at step $n$ is $\tau_n = \sum_{i=1}^{n} \Delta_i$, where $\Delta_i = F_i$ is the Fibonacci integration latency.
\end{definition}

\begin{definition}[Physical Time]
\label{def:phys-time}
\emph{Physical time} $t$ is the continuous geometric parameter arising from the Kinematic-Dynamic Equivalence (Lemma~5.10 of~\cite{pen}):
\begin{equation}
    \nextmod X \;\simeq\; X^{\D}.
\end{equation}
Physical time exists only after Step~15, when the ``next'' modality $\nextmod$ is identified with the tangent endofunctor $(-)^{\D}$, converting discrete temporal succession into continuous infinitesimal displacement.
\end{definition}

The two clocks have entirely different character.
Algorithmic time is discrete, external to the library, and measured in structural ticks.
Physical time is continuous, internal to the DCT, and measured in infinitesimal tangent displacements.
There is no meaningful sense in which $\tau_{14}$ is ``earlier than'' $t = 0$; the two quantities inhabit different ontological categories, related only by the phase transition at Step~15 that converts one into the other.

\begin{remark}[Why ``before the Big Bang'' is a category error]
\label{rem:category-error}
Asking ``what happened before the Big Bang'' presupposes that physical time extends to the left of $t = 0$.
In PEN, the pre-Step~15 epoch is not a prior interval of $t$ but a structurally distinct generative process indexed by $\tau$.
The question dissolves: nothing happened \emph{before} $t = 0$ in physical time, because physical time does not exist before Step~15.
What happened \emph{instead} of physical time was algorithmic assembly---a categorically different process that produced the conditions under which physical time could exist at all.
\end{remark}


% ============================================
% 3. STEP 15 IGNITION
% ============================================
\section{Step~15 Ignition: The Birth of Physical Time \texorpdfstring{($t = 0$)}{(t = 0)}}
\label{sec:ignition}

\subsection{The Arrival of LTL}
\label{sec:ltl}

The Step~15 temporal-cohesive shell extends the library $\B_{14}$ with four structural components~\cite{pen}:

\begin{enumerate}[label=(\arabic*), nosep]
    \item \textbf{Spatial Logic:} The cohesive interface already present in $\B_{14}$, namely $(\flat \dashv \sharp,\; \Pi_{\mathrm{coh}} \dashv \Disc)$.
    \item \textbf{Temporal Core:} The modalities $\nextmod$ (``next'') and $\eventually$ (``eventually'') from Linear Temporal Logic, following Nakano~\cite{nakano}.
    \item \textbf{Spatial-Temporal Exchange:} Explicit commuting laws between temporal and cohesive structure, exemplified by $\nextmod(\flat X) \simeq \flat(\nextmod X)$ and $\sharp(\eventually X) \simeq \eventually(\sharp X)$.
    \item \textbf{Guarded Coherence:} The temporal self-interaction and elimination clauses that make the shell executable as a typed AST.
\end{enumerate}

\noindent The specification cost is $\kappa = 8$, audited as the eight temporal-cohesive clauses of the selected AST~\cite{pen}.
In the current strict machine canon, the generative capacity is $\nu = 88$, yielding $\rho = 11.00$ and clearing the Step~15 bar of $7.51$ by the largest absolute overshoot in the sequence~\cite{pen}.
The infinitesimal object $\D$ belongs to the semantic completion of this shell rather than to the primitive Step~15 syntax.

\subsection{The Kinematic-Dynamic Equivalence as Geometrization of Time}
\label{sec:kde}

The central mechanism is the Kinematic-Dynamic Equivalence: in the semantic completion of the temporal-cohesive shell, the discrete temporal modality $\nextmod$ is naturally equivalent to the continuous tangent endofunctor $(-)^{\D}$~\cite{pen}.

This is not a primitive stipulation of the selected AST but a \emph{forced semantic identification}.
The exchange laws with the cohesive modalities and the guarded coherence clauses constrain $\nextmod$ so strongly that, in the DCT reading, it is represented by tangent exponentiation.
The proof does not begin by postulating $\D$ in the syntax; rather, $\D$ is derived as the infinitesimal shadow of the first-order temporal-cohesive endofunctor selected by the engine~\cite{pen}.

The physical content of this equivalence is profound: \emph{the passage of time is the same operation as infinitesimal spatial displacement}.
What the temporal logic calls ``the next state'' is what synthetic differential geometry calls ``an infinitesimal tangent vector.''
This identification establishes $t = 0$---not as a point on a pre-existing timeline, but as the moment at which the concept of temporal succession acquires geometric content.

\subsection{The Guarded Corecursion Engine}
\label{sec:corecursion}

The temporal logic of Nakano~\cite{nakano} carries a powerful recursion principle: L\"ob's guarded fixpoint rule,
\begin{equation}
\label{eq:lob}
    \fix : (\nextmod A \to A) \to A.
\end{equation}
This rule states: given a function that can produce an $A$ provided it has access to ``tomorrow's'' $A$, one can construct a definite $A$ today.
The guardedness condition (the $\nextmod$ on the left of the arrow) prevents circular self-reference by ensuring that present computation may depend on future values only through the temporal delay.

Applied to the DCT's kinematic library, L\"ob's rule has the following operational effect.
A discrete-time dynamical system is a map $f : X \to \nextmod X$---given a current state, compute the next state.
Under the Kinematic-Dynamic Equivalence, this becomes:
\begin{equation}
\label{eq:vf}
    f : X \to X^{\D},
\end{equation}
which is precisely a \emph{vector field} in synthetic differential geometry: a section of the tangent bundle assigning an infinitesimal displacement to each point.
The guarded corecursion principle then unfolds any such vector field into a global trajectory:
\begin{equation}
\label{eq:unfold}
    \unfold_f : X \;\to\; \Stream\, X.
\end{equation}
This is the cosmological spark.
The 14~static structures of $\B_{14}$ are the fuel; $\nextmod$ is the ignition mechanism; L\"ob's rule is the engine that converts a local infinitesimal prescription into perpetual, autonomous motion.
Before Step~15, the library describes spaces.
After Step~15, the library computes trajectories through those spaces.

\begin{proposition}[Step~15 as the onset of dynamics]
\label{prop:onset}
Prior to the realization of the DCT, no structure in $\B_{14}$ can express temporal evolution: there is no type former producing streams, no fixpoint principle, and no identification of succession with tangent displacement.
The DCT is the \emph{unique minimal} extension of $\B_{14}$ that introduces autonomous dynamics, and it does so with $\kappa = 8$ clauses.
\end{proposition}

\begin{proof}
Dynamics requires (i)~a notion of succession, (ii)~a fixpoint or corecursion principle to unfold finite prescriptions into infinite trajectories, and (iii)~compatibility with the existing geometric library so that trajectories inhabit the already-constructed spaces.
Requirements (i) and (ii) are supplied by $\nextmod$ and L\"ob's rule; requirement~(iii) is enforced by the temporal-cohesive exchange laws and guarded coherence clauses.
No proper subset of these eight clauses suffices: removing any temporal modality eliminates succession; removing the exchange laws severs temporal evolution from the geometric library; removing guarded coherence destroys the executable dynamics.
The infinitesimal interpretation is then recovered semantically from this shell rather than paid for as a primitive clause.
\end{proof}


% ============================================
% 4. COSMIC INFLATION
% ============================================
\section{Structural Inflation and the Preconditions for Cosmological Inflation}
\label{sec:inflation}

\subsection{\texorpdfstring{The $\rho = 11.00$ Efficiency Peak}{The rho = 11.00 Efficiency Peak}}
\label{sec:rho-spike}

The strict machine-verified efficiency ratio at Step~15 is $\rho = 88/8 = 11.00$, clearing the bar ($7.51$) by an absolute overshoot of $3.49$---the largest in the entire 15-step sequence~\cite{pen}.
This ``efficiency singularity'' is the PEN counterpart of the cosmological singularity, but with a crucial difference: the mathematics does not break down.
Where General Relativity produces divergent curvature scalars at $t = 0$, PEN produces a finite, well-defined maximum of algorithmic efficiency.
The singularity is not a failure of the formalism but a signature of \emph{optimal synthesis}---the point at which the library's combinatorial capacity is most powerfully leveraged.

\begin{remark}[Efficiency singularity vs.\ curvature singularity]
\label{rem:singularity-contrast}
In GR, the initial singularity signals the breakdown of the theory: physical quantities (density, temperature, curvature) diverge, and the Einstein equations lose predictive power~\cite{hawking-ellis,wald}.
In PEN, the ``singularity'' at Step~15 is a computable finite maximum of a well-defined ratio.
It arises not from a failure of description but from the temporal-cohesive shell's measured interaction payload with the existing library.
The manuscript now treats the strict machine total $\nu = 88$ as canonical and no longer promotes the older structural-AST tally to theorem-level status~\cite{pen}.
The analogy is precise: both mark the moment of cosmological origin.
The disanalogy is equally precise: one is a breakdown, the other is a peak of algorithmic performance.
\end{remark}

\subsection{\texorpdfstring{The $\nu = 88$ Explosion and the Inflationary Epoch}{The nu = 88 Explosion and the Inflationary Epoch}}
\label{sec:nu-explosion}

Cosmic inflation is the hypothesis that the very early universe underwent a brief period of exponentially rapid expansion of the metric scale factor, during which the number of accessible physical degrees of freedom increased by many orders of magnitude~\cite{guth,linde}.
In PEN, the analogous phenomenon is not the spatial expansion itself but the combinatorial schema explosion at Step~15 that first makes such a phase dynamically expressible.
The claim of this section is therefore one of \emph{structural inflation}: the sudden release of executable geometric and temporal degrees of freedom.
It does \emph{not} identify the logical count $\nu = 88$ with the metric scale factor $a(t)$.

The mechanism is the Combinatorial Schema Synthesis theorem~\cite{pen}: when the DCT's temporal modalities ($\nextmod$, $\eventually$) are introduced into a library of 14~structures already equipped with spatial modalities ($\flat, \sharp, \Disc, \Pi_{\mathrm{coh}}$), the cross-interaction at depth~2 produces a superlinear expansion of new inference schemas.
In the current strict lane, this temporal-cohesive shell is evaluated at $\kappa = 8$ and yields the canonical machine value $\nu = 88$.
Representative schema families include:

\begin{itemize}[nosep]
    \item Guarded temporal endomorphism and recursion patterns built from $\nextmod$ and $\eventually$.
    \item Spatial-temporal exchange laws that tie the new temporal operators to the pre-existing cohesive modalities.
    \item Transport and flow-like interaction with the already selected cohesive and differential layers.
\end{itemize}

The picture is one of sudden, explosive logical cross-interaction: 14~structures that previously coexisted as independent library entries are now woven into a single dynamical fabric, with each pair generating new eliminators, new homotopy schemas, and new compositional types.
This is the algorithmic analogue of the release of latent degrees of freedom in an inflationary phase transition: the physical ingredients are not injected from outside but are \emph{unlocked} by the synthesis of previously disconnected structure.

\begin{proposition}[Structural inflation as a necessary precondition]
\label{prop:structural-inflation}
The combinatorial amplification at Step~15 ($\nu = 88$, compared to $\nu \leq 63$ at all prior steps) is a necessary structural precondition for any later physical inflationary dynamics.
It is not itself the claim that the metric scale factor $a(t)$ has already undergone exponential expansion.
Rather, Step~15 is the first stage at which the ingredients required to formulate such dynamics coexist inside a single executable framework: differentiable manifolds, Lorentzian metric structure, temporal evolution, and guardedly unfoldable fields.
\end{proposition}

\begin{proof}[Proof sketch]
Before Step~15, the relevant geometric ingredients are present only as static library entries in $\B_{14}$; after Step~15, they are bound together by the temporal-cohesive shell and acquire dynamical meaning through the Kinematic-Dynamic Equivalence.
A Guth-Linde style inflationary model requires not merely a manifold and a scalar field, but an equation of motion propagating that field on a spacetime background.
PEN therefore identifies the earliest structural point at which physical inflation can even be written down.
Which concrete inflaton potential or low-energy effective action is realized is a separate phenomenological question, deferred to~\cite{gr-companion}.
\end{proof}


% ============================================
% 5. THE ARROW OF TIME
% ============================================
\section{The Arrow of Time and Minimal Entropy}
\label{sec:arrow}

\subsection{The Causal Arrow from Guardedness}
\label{sec:causal-arrow}

Why does time flow in one direction?
In thermodynamic accounts, the arrow of time is derived from the second law and the assumption of a low-entropy initial condition---but this merely pushes the question back to why the initial condition was special.
In PEN, the arrow of time is not a thermodynamic consequence but a \emph{type-theoretic primitive}: it is built into the logical structure of the $\nextmod$ operator.

The guardedness condition in L\"ob's rule~\eqref{eq:lob} is asymmetric by construction: a term of type $\nextmod A \to A$ can use tomorrow's value to compute today's value, but there is no dual rule $A \to \nextmod A$ that would allow today's value to determine tomorrow's independently of the dynamical law.
More precisely, the only way to produce a term of type $\nextmod A$ is by applying a dynamical system $f : A \to \nextmod A$ to a current state; one cannot extract a future state except by running the system forward.

\begin{theorem}[Type-theoretic irreversibility]
\label{thm:irreversibility}
In the DCT, the guarded corecursion principle $\fix : (\nextmod A \to A) \to A$ is not invertible: there is no general term of type $A \to (\nextmod A \to A)$ that recovers the generating function from its fixpoint.
Consequently, the temporal unfolding $\unfold_f : X \to \Stream\, X$ is causally one-directional: given a stream, one can extract any finite prefix by iterated application of $\nextmod$, but one cannot in general recover the state at tick $n$ from the state at tick $n + k$ for $k > 0$.
\end{theorem}

\begin{proof}[Proof sketch]
Suppose for contradiction that a general inverse $\fix^{-1} : A \to (\nextmod A \to A)$ existed.
Then for any type $A$, one could construct a term of type $\nextmod A \to A$ from any element of $A$, which when fed back into $\fix$ would yield a fixpoint independent of the dynamical law $f$.
This violates the productivity guarantee of guarded corecursion: the $\nextmod$ guard ensures that each step of the unfolding makes genuine computational progress.
A general inverse would collapse the guarded modality, making $\nextmod A \simeq A$, which contradicts the non-triviality of temporal logic in the DCT (the temporal-cohesive exchange laws require $\nextmod$ to remain distinguishable from the identity).
\end{proof}

\begin{remark}[Execution asymmetry vs.\ geometric reversibility]
\label{rem:reversibility}
The asymmetry proved in \Cref{thm:irreversibility} concerns the production of the runtime trace, not the symmetry group of every internal law of motion.
Once Step~15 has supplied a geometric vector field $v : \U \to \U^{\D}$, that field may integrate to a reversible diffeomorphism, a symplectic flow, or an equation admitting a formal time-reversal involution.
One may therefore solve a microscopic law backward inside the already-generated geometry.
What one cannot do is \emph{un-run} the universe's guarded execution to recover the present by retrocausal compilation.
The runtime computes forward even when some of the equations it hosts are geometrically reversible.
\end{remark}

\begin{corollary}[Type safety excludes closed timelike curves]
\label{cor:ctc}
In the guarded fragment of the DCT, any putative closed timelike curve requiring a state to be fed into its own causal past would demand discharge of a $\nextmod$-guard without forward progress.
No such term is typable.
Hence closed timelike curves are excluded at the level of the universe's native type safety.
\end{corollary}

\begin{proof}
A CTC would amount to a circular dependency in which a future state is extracted from a stream and reintroduced as an input to its own earlier production.
By \Cref{thm:irreversibility}, the guarded unfolding map is not generally invertible, and by guarded typing rules the delay on $\nextmod$ cannot be eliminated except through forward execution.
Therefore the required self-feeding term is ill-typed.
\end{proof}

\subsection{The Past Hypothesis as Algorithmic Compression}
\label{sec:past-hypothesis}

The ``Past Hypothesis''---the empirical observation that the universe began in a state of extraordinarily low entropy~\cite{penrose-road,carroll}---is one of the deepest unexplained facts in physics.
Statistical mechanics can explain why entropy increases \emph{given} a low-entropy initial state, but it cannot explain why the initial state was special.
The low-entropy boundary condition is typically imposed by hand as a cosmological axiom.

PEN provides a constructive explanation.
The $t = 0$ state---the state of the universe at the moment of runtime ignition---is not an arbitrary initial condition.
It is the \emph{output} of 15~steps of strict efficiency maximization under Kolmogorov-complexity-minimal specification.

\begin{theorem}[Minimal realized description at $t = 0$]
\label{thm:past-hypothesis}
The physical state at the onset of dynamics ($t = 0$) is specified by the complete DCT library $\B_{15}$, whose total specification complexity is $\kappa_{\mathrm{total}} = \sum_{i=1}^{15} \kappa_i = 92$ clauses (equivalently, 610~bits in MBTT encoding~\cite{pen}).
Within the PEN selection dynamics, this is the shortest realized description compatible with the full kinematic framework.
Any alternative bar-clearing trajectory obeying the same local minimal-overshoot rule and producing equivalent kinematic capacity requires equal or greater total $\kappa$.
\end{theorem}

\begin{proof}[Proof sketch]
The selection dynamics at each step chooses the candidate with minimal overshoot (Axiom~4 of~\cite{pen}), with ties broken by minimal $\kappa$.
Within that admissible class of bar-clearing trajectories, the realized path is therefore greedily minimal at every stage.
Any alternative trajectory that reaches equivalent kinematic capacity while obeying the same acceptance rule must match or exceed the accumulated specification cost of the realized path.
\end{proof}

\begin{corollary}[Past Hypothesis in Zurek's algorithmic sense]
\label{cor:zurek-past}
Zurek's account of physical entropy assigns to a definite microstate an algorithmic contribution given by the length of its shortest specification, together with the usual coarse-grained missing information~\cite{zurek-complexity}.
Since PEN realizes the ignition state by the shortest description in its admissible class, the algorithmic contribution to physical entropy is minimized at $t = 0$.
If ignition is also taken to precede any substantial coarse-grained branching of macrostates, then the total physical entropy is correspondingly minimal there, yielding a constructive realization of the Past Hypothesis.
\end{corollary}

\begin{proof}
The first clause is Zurek's decomposition of physical entropy.
The second follows from \Cref{thm:past-hypothesis}: the ignition state is generated by the shortest realized description compatible with the PEN constraints.
The final sentence is the cosmological interpretation obtained by applying that decomposition at the boundary between compilation and runtime.
\end{proof}

\begin{remark}[What the 610-bit claim does and does not mean]
\label{rem:entropy}
The 610-bit figure is a statement about the shortest description of the ignition state, not a literal calculation of the late universe's thermodynamic entropy.
Its force lies in fixing the algorithmic part of the entropy budget at a minimum compatible with PEN.
Once runtime begins, coarse-grained multiplicity can grow rapidly, recovering the familiar thermodynamic arrow without sacrificing the compressed origin.
In this precise sense, PEN turns the Past Hypothesis from a brute boundary condition into the output of a compression principle.
\end{remark}


% ============================================
% 6. THE UNIVALENT HORIZON AS COSMOLOGICAL FIREWALL
% ============================================
\section{The Univalent Horizon as the Cosmological Firewall}
\label{sec:firewall}

\subsection{Resolving the Halting Problem: Structural Termination}
\label{sec:termination}

The PEN Generative Sequence does not ``die'' at Step~15.
It halts---in the precise sense of Theorem~5.12 of~\cite{pen}---because the external generator is rendered obsolete.

The argument has three steps:
\begin{enumerate}[nosep]
    \item The DCT internalizes the evolutionary mechanism as a geometric flow on $\U$ (Lemma~5.11 of~\cite{pen}).
    The metamathematical process of ``searching for the next type'' is representable as a vector field $v : \U \to \U^{\D}$, unfolded into a trajectory by guarded corecursion.
    \item By univalence, paths in $\U$ are type equivalences: $(A =_{\U} B) \simeq (A \simeq B)$.
    Integrating the vector field produces a continuous family of types connected by paths in $\U$, hence a family of pairwise equivalent types.
    The flow is permanently confined to the connected component of $\U$ containing $\B_{15}$.
    \item Any candidate $Y$ proposed at Step~16 is either (a)~internally derivable (contributing $\nu = 0$) or (b)~structurally alien to the $d = 2$ Coherence Window (incurring $\Delta_{16} = F_{16} = 987$ integration latency, making $\rho \ll \mathrm{Bar}_{16} \approx 8.80$).
    In either case, the bar is permanently uncleared.
\end{enumerate}

The universe does not stop evolving at Step~15; it stops \emph{needing external structural input} at Step~15.
The distinction is between a system that cannot continue (failure) and a system that has achieved self-sufficiency (closure).
The Generative Sequence halts because the DCT is a fixpoint of the selection dynamics, not because the dynamics has run out of candidates.

\subsection{Epistemic Isolation: The Planck Boundary Analogy}
\label{sec:epistemic}

The closest standard cosmological analogue of the Univalent Horizon is not the Cosmic Microwave Background at recombination ($t \approx 380{,}000$ years), which is only an electromagnetic opacity surface, but the Planck boundary ($t \lesssim 10^{-43}\,\mathrm{s}$), where classical spacetime is no longer expected to furnish a reliable description~\cite{mukhanov}.
PEN provides a structural analogue of that deeper boundary: not merely a limit on photon propagation, but a limit on what an internal observer can represent as a physical observable.

\begin{proposition}[Epistemic isolation of the compilation epoch]
\label{prop:epistemic}
Internal observers---agents whose perceptual and computational apparatus is itself a structure within $\B_{15}$---cannot construct an observable that probes the discrete algorithmic steps (1--14) that built the topos.
The pre-Bang compilation epoch is accessible only via type-theoretic deduction, not physical measurement.
\end{proposition}

\begin{proof}[Proof sketch]
An ``observation'' in the DCT is a continuous geometric flow: a section of some bundle, a functional on a Hilbert space, or a path in a moduli space.
All such flows are terms in the internal language of the DCT, and by the Univalent Horizon (Theorem~5.12 of~\cite{pen}), they are confined to the connected component of $\U$ containing $\B_{15}$.
The discrete algorithmic steps $\tau_1, \ldots, \tau_{14}$ are \emph{not} connected to $\B_{15}$ by paths in $\U$---they belong to the external metatheory, not the internal geometry.
Any attempt to construct a physical observable that ``sees'' the compilation epoch would require a path crossing the Univalent Horizon, which is precisely what the horizon forbids.
\end{proof}

\begin{remark}[The Planck boundary as a physical shadow of the Univalent Horizon]
\label{rem:planck}
The Planck boundary is the stage at which continuum spacetime ceases to be a trustworthy effective description and pre-geometric or quantum-gravitational structure is expected to dominate~\cite{mukhanov}.
In PEN, the Univalent Horizon plays the analogous structural role.
The limitation is not just that a particular signal channel thermalizes; it is that internal continuous observables do not extend across the compilation/runtime boundary at all.
The analogy is therefore to a breakdown of representability, not merely to an opacity of electromagnetic matter.
\end{remark}


% ============================================
% 7. CONCLUSION
% ============================================
\section{Conclusion}
\label{sec:conclusion}

The standard approach to cosmogenesis looks for a physical, kinetic mechanism that \emph{triggered} the Big Bang---a quantum fluctuation, a tunneling event, a collision of branes.
Each such proposal presupposes the very structures (quantum fields, spacetime, dynamics) whose origin it is supposed to explain.

PEN inverts the explanatory order.
The Big Bang is not an event that happened within physics; it is the mathematically necessary boundary at which physics \emph{begins}.
Specifically:

\begin{enumerate}[nosep]
    \item \textbf{The compilation epoch} (Steps~1--14) assembles the kinematic framework---types, spaces, cohesion, metrics, Hilbert functionals---in discrete algorithmic time $\tau$, without presupposing continuous dynamics.
    \item \textbf{Runtime ignition} (Step~15) introduces the DCT, whose Kinematic-Dynamic Equivalence ($\nextmod X \simeq X^{\D}$) geometrizes time, converting discrete temporal logic into continuous infinitesimal structure.
    L\"ob's guarded corecursion sets the static library into perpetual, autonomous motion.
    \item \textbf{Structural inflation} is the combinatorial schema explosion ($\nu = 88$) produced by the sudden cross-interaction of 14~previously disconnected kinematic structures; it supplies the precondition for, rather than a replacement of, any later metric inflation.
    \item \textbf{The arrow of time} is the type-theoretic asymmetry of guarded execution: the runtime computes forward even when some internal microscopic laws remain geometrically time-reversal symmetric.
    \item \textbf{The Past Hypothesis} is realized as the minimized algorithmic contribution to physical entropy, in Zurek's sense, at the moment the runtime first ignites.
    \item \textbf{The Univalent Horizon} is a Planck-like cosmological firewall separating the pre-physical compilation epoch from the self-computing physical runtime.
    The pre-Bang universe is epistemically isolated from internal observers, accessible only through the type-theoretic deduction that PEN provides.
\end{enumerate}

On the structural-realist reading adopted here, the physical universe is the runtime execution of the Dynamical Cohesive Topos.
We must stop looking for a first cause within physics and instead recognize the Big Bang as the halting condition of the universe's foundational compilation phase---the moment at which the formal system achieves self-sufficiency and begins computing its own future.


% ============================================
% REFERENCES
% ============================================
\begin{thebibliography}{99}

\bibitem{pen}
H.S.\ Lande, ``The Principle of Efficient Novelty: The Algorithmic Origin of the Kinematic Framework of Physics,'' \texttt{pen\_unified.tex}, 2026.

\bibitem{gr-companion}
H.S.\ Lande, ``Structural Emergence of 4D Lorentzian Geometry from the Principle of Efficient Novelty,'' companion paper, 2026.

\bibitem{hawking-ellis}
S.W.\ Hawking and G.F.R.\ Ellis, \emph{The Large Scale Structure of Space-Time}, Cambridge University Press, 1973.

\bibitem{wald}
R.M.\ Wald, \emph{General Relativity}, University of Chicago Press, 1984.

\bibitem{nakano}
H.\ Nakano, ``A Modality for Recursion,'' \emph{Proceedings of the 15th Annual IEEE Symposium on Logic in Computer Science (LICS)}, 2000.

\bibitem{guth}
A.H.\ Guth, ``Inflationary Universe: A Possible Solution to the Horizon and Flatness Problems,'' \emph{Physical Review D} 23 (1981), 347--356.

\bibitem{linde}
A.D.\ Linde, ``A New Inflationary Universe Scenario: A Possible Solution of the Horizon, Flatness, Homogeneity, Isotropy, and Primordial Monopole Problems,'' \emph{Physics Letters B} 108 (1982), 389--393.

\bibitem{mukhanov}
V.F.\ Mukhanov, \emph{Physical Foundations of Cosmology}, Cambridge University Press, 2005.

\bibitem{penrose-road}
R.\ Penrose, \emph{The Road to Reality}, Jonathan Cape, 2004.

\bibitem{carroll}
S.M.\ Carroll, ``The Cosmological Constant,'' \emph{Living Reviews in Relativity} 4 (2001), 1.

\bibitem{zurek-complexity}
W.H.\ Zurek, ``Algorithmic Randomness and Physical Entropy,'' \emph{Physical Review A} 40 (1989), 4731--4751.

\bibitem{hott}
The Univalent Foundations Program, ``Homotopy Type Theory: Univalent Foundations of Mathematics,'' Institute for Advanced Study, 2013.

\end{thebibliography}

\end{document}
